\documentclass[11pt]{article}
\usepackage{amsmath,amssymb}
\usepackage{geometry}
\usepackage{booktabs}
\usepackage{hyperref}
\geometry{margin=1in}

\newcommand{\grad}{\nabla}
\newcommand{\R}{\mathbb{R}}

\title{Multi-Label Potential Field Reconstruction on the BCC Lattice:\\
Energy Terms and the Jacobi Solve}
\author{}
\date{}

\begin{document}
\maketitle

\section{Goals}

The lattice carries, at every node, a discrete material \emph{label} and a
real-valued \emph{potential}. For each label $\ell$ present in the domain we
want the collection of node potentials, reinterpreted as a scalar field, to
behave like a well-formed distance-like function: positive and large deep
inside label $\ell$'s territory, crossing zero at the interface with any
other label, and doing so smoothly enough that a multi-material surface
(including non-manifold triple-junctions where three labels meet along a
curve) can be extracted directly from it.

Four properties are wanted of this field, in decreasing order of how
fundamental they are to making the others meaningful:

\begin{enumerate}
  \item \textbf{Eikonal (unit-gradient) shaping.} Wherever a node is
  genuinely responsible for representing the local edge of its label's
  territory, the field's gradient magnitude should be close to $1$ --
  neither collapsing to a flat, informationless field nor blowing up.
  \item \textbf{Smoothing.} Nearby same-label potentials should vary
  gradually; the field should not carry high-frequency noise unrelated to
  actual geometric structure.
  \item \textbf{Juncture straightness.} Where three labels meet, the local
  triple-junction curve should be locally straight -- consecutive
  tetrahedra spanning a junction should agree on its direction, not kink.
  \item \textbf{Connectivity / volume preservation.} Each label's
  reconstructed region should retain a sensible target volume (and,
  eventually, topological connectivity) rather than drifting or shrinking
  away purely as a side effect of satisfying the other three terms.
\end{enumerate}

Item 1 is listed first deliberately: it is a \emph{prerequisite} for items
2--4 rather than a peer among equals. Every other term's residual is
constructed from gradients of the same underlying per-label field, and
every one of those residuals is trivially -- and uselessly -- satisfied by
letting those gradients collapse to zero. The Eikonal term is the only term
whose energy actively resists that collapse; without it (or with it too
weak relative to the others), the solve is free to flatten the field
entirely, at which point smoothing, junction, and connectivity residuals
are all vacuously zero and carry no information about the actual geometry.
Section~\ref{sec:eikonal} is therefore treated first among the four
mathematical sections below, both because of this dependency and because
its stencil and derivative machinery (per-tet affine gradient
reconstruction) is reused verbatim by the juncture-straightness term.

\section{Field Representation}

\subsection{Lattice}

Nodes live on a body-centered-cubic (BCC) lattice, split into two
sublattices: $A$-nodes (the underlying simple-cubic grid's corners, one per
integer triple $(i,j,k)$, $0 \le i,j,k < N$ for grid resolution $N$) and
$B$-nodes (the cube centers, $(i,j,k)$ with $0 \le i,j,k < N-1$). Every node
carries exactly one (label, potential) pair, $(\text{label}(x), \phi(x))$
with $\phi(x) \ge 0$.

Both sublattices are addressed through a shared change of basis into a
virtual simple-cubic \emph{$q$-space}: $q = (i+j,\,i+k,\,j+k)$ for an
$A$-node and $q=(i+j+1,\,i+k+1,\,j+k+1)$ for a $B$-node, with node type
recovered as the parity of $q_x+q_y+q_z$ (even $=A$, odd $=B$). Every unit
cube of $q$-space decomposes into $6$ tetrahedra sharing the cube's main
diagonal; every consumer below enumerates tetrahedra, a node's incident
tetrahedra, or a tetrahedron's face-adjacent partner purely from a fixed,
small set of corner-offset tables and a tet or node index -- there is no
explicit connectivity buffer.

\subsection{Per-label confidence field}

For a node $x$ and a label $\ell$, define the signed \emph{confidence
field}
\begin{equation}
  G_\ell(x) \;=\; \phi(x)\cdot\sigma(x,\ell), \qquad
  \sigma(x,\ell) = \begin{cases} +1 & \text{label}(x) = \ell \\ -1 &
  \text{label}(x) \ne \ell \end{cases}.
  \label{eq:confidence}
\end{equation}

Every node contributes its own real potential magnitude to \emph{every}
label's field, sign-flipped by agreement -- this is not a same-label
gate or exclusion. $G_\ell$ is positive where $\ell$ locally dominates and
negative elsewhere, and its zero level set is exactly the boundary of
label $\ell$'s territory. Equation~\eqref{eq:confidence} is the single
construction reused by every term below.

Within one tetrahedron $\tau$ with world-space corner positions
$P_0,\dots,P_3$, any corner-affine scalar field with values $g_0,\dots,g_3$
extends uniquely to an affine function of position, with constant gradient
\begin{equation}
  \grad g \;=\; \sum_{c=0}^{3} g_c\, w_c(\tau),
  \label{eq:shapegrad}
\end{equation}
where $w_c(\tau) \in \R^3$ are the tetrahedron's shape-function gradients
(the reciprocal/dual basis of its edge frame):
\begin{equation}
  e_1 = P_1{-}P_0,\quad e_2 = P_2{-}P_0,\quad e_3 = P_3{-}P_0,\qquad
  V = e_1\cdot(e_2\times e_3),
\end{equation}
\begin{equation}
  w_1 = \frac{e_2\times e_3}{V},\quad
  w_2 = \frac{e_3\times e_1}{V},\quad
  w_3 = \frac{e_1\times e_2}{V},\quad
  w_0 = -(w_1+w_2+w_3).
  \label{eq:dualbasis}
\end{equation}

Applying \eqref{eq:shapegrad} with $g_c = G_\ell(c)$ for $\tau$'s four
corners gives the per-tet, per-label gradient $\grad G_\ell|_\tau$ used
throughout Sections~\ref{sec:eikonal} and \ref{sec:junction}.

\section{The Jacobi Scheme as a Solve Tool}
\label{sec:jacobi}

All four terms are folded into one shared local update, evaluated
independently and simultaneously at every node each sweep (Jacobi, not
Gauss--Seidel: every neighboring unknown is treated as frozen at its
value from the start of the sweep).

The total energy is a weighted sum over terms,
\begin{equation}
  E[\phi] \;=\; \sum_{k} w_k\, E_k[\phi],
  \qquad k \in \{\text{Eikonal},\ \text{Smoothing},\ \text{Junction},\
  \text{Volume}\},
\end{equation}
and each term's energy is (or, for Volume, will be) a sum of squared scalar
residuals over some geometric stencil of tet or node tuples $s$,
\begin{equation}
  E_k[\phi] \;=\; \sum_{s} \rho_{k,s}\, r_{k,s}[\phi]^2,
\end{equation}
with $\rho_{k,s} \ge 0$ an optional, $\phi$-independent per-instance weight
(e.g.\ a gate or a double-counting correction).

For a single scalar unknown $\phi_x$ (one node's own potential) being
updated this sweep, every term contributes a local gradient and a
Gauss--Newton curvature (diagonal) estimate:
\begin{equation}
  \frac{\partial E_k}{\partial \phi_x} \;=\; \sum_s \rho_{k,s}\cdot 2\,
  r_{k,s}\,\frac{\partial r_{k,s}}{\partial \phi_x},
  \qquad
  D_k(x) \;:=\; \sum_s \rho_{k,s}\cdot 2\left(\frac{\partial
  r_{k,s}}{\partial \phi_x}\right)^{\!2}.
  \label{eq:gaussnewton}
\end{equation}
$D_k(x)$ discards the residual's own second-derivative term, which is what
keeps it a sum of squares and hence \emph{always non-negative} -- essential
for a stable descent direction when every neighboring unknown is being
updated simultaneously rather than sequentially. This is why every term
below is deliberately written in Gauss--Newton form even where a term is
nonlinear enough that a different, exact linearization would exist.

The combined update is
\begin{equation}
  \text{grad}(x) = \sum_k w_k\,\frac{\partial E_k}{\partial\phi_x}, \qquad
  \text{diag}(x) = \sum_k w_k\, D_k(x),
\end{equation}
\begin{equation}
  \Delta\phi_x \;=\; \mathrm{clamp}\!\left(
    \frac{-\,\text{grad}(x)}{\text{diag}(x) + \varepsilon},\
    -\Delta_{\max},\ \Delta_{\max}
  \right), \qquad \phi_x \leftarrow \phi_x + \Delta\phi_x,
  \label{eq:update}
\end{equation}
with $\varepsilon$ a fixed diagonal floor (preventing division blow-up when
every term's diagonal contribution happens to vanish at a node) and
$\Delta_{\max}$ a hard per-sweep trust-region clamp. The clamp is required
because a Jacobi step neglects all off-diagonal coupling by construction;
without it, the coupled nonlinear system's naive per-node update
overshoots and diverges over successive sweeps.

If a node's label is not held fixed (an unlabeled/free node, as opposed to
a ground-truth node), a relabeling decision is made \emph{after} the
potential update, whenever $\phi_x + \Delta\phi_x$ would go negative: the
node adopts whichever neighboring label locally dominates and its
potential is reset to a small positive head start, so that it participates
in the next sweep's terms on comparable footing with its already-settled
neighbors rather than needing many sweeps of $\Delta_{\max}$-limited growth
to catch up.

\section{Eikonal Floor Term}
\label{sec:eikonal}

\subsection{Purpose}

This term keeps $|\grad G_{\text{label}(x)}|$ near $1$ wherever node $x$ is
genuinely the local peak of its own label's confidence field -- i.e.\
wherever $x$ is responsible for representing the boundary there. It is a
one-sided \emph{floor}: it only pushes the gradient magnitude \emph{up}
when it is below $1$, and only where $x$ is that local peak; it does not
penalize a large or already-unit gradient, and it does not act at all at a
node that is not locally the maximum (that responsibility instead falls,
symmetrically, to whichever corner actually is the local maximum, when
\emph{that} node is processed as the target).

\subsection{Stencil}

For target node $x$ with label $\ell = \text{label}(x)$, let $\tau$ range
over $x$'s $24$ ring-1 incident tetrahedra (every tetrahedron that has $x$
as one of its four corners). For each such $\tau$ with corners $c=0,\dots,3$
and $x$ at local index $m$:
\begin{equation}
  G_c \;=\; G_\ell(c) \quad\text{(Eq.~\eqref{eq:confidence}, evaluated in
  label $\ell$ for every corner of $\tau$)}.
\end{equation}

\subsection{Gating: strict local maximum, softened}

Define the margin by which $x$ exceeds the best of the other three corners,
\begin{equation}
  \mathrm{margin} \;=\; G_m - \max_{c \ne m} G_c,
\end{equation}
and a gate
\begin{equation}
  \mathrm{gate} \;=\; \mathrm{saturate}\!\left(\frac{\mathrm{margin}}
  {\epsilon_{\text{floor}}}\right) \in [0,1], \qquad \epsilon_{\text{floor}}
  = 10^{-4}.
  \label{eq:gate}
\end{equation}
A hard $0/1$ strict-maximum test would flip sweep-to-sweep near a near-tie
between two corners; \eqref{eq:gate} is a narrow anti-chatter ramp (scaled
to the same small constant used as a diagonal floor elsewhere, not a wide
physical smoothing radius) that only smooths the immediate tie-breaking
boundary. Terms with $\mathrm{gate}=0$ are skipped entirely.

\subsection{Residual and fold-in}

Reconstruct the real per-tet gradient via \eqref{eq:shapegrad},
\begin{equation}
  \grad G \;=\; \sum_{c=0}^{3} G_c\, w_c(\tau), \qquad
  |\grad G| \text{ its magnitude.}
\end{equation}
If $|\grad G| < 1$, the one-sided residual is
\begin{equation}
  r \;=\; 1 - |\grad G|,
\end{equation}
contributing energy $\mathrm{gate}\cdot r^2$ (zero, with zero derivative,
once $|\grad G|\ge 1$ -- a $C^1$-continuous hinge). Since $x$'s own corner
always agrees with its own label, $\partial G_m/\partial\phi_x = +1$
identically, so $\partial(\grad G)/\partial\phi_x = w_m(\tau)$ with no sign
branch, and
\begin{equation}
  \frac{\partial |\grad G|}{\partial \phi_x} = \frac{\grad G \cdot
  w_m(\tau)}{\max(|\grad G|,\,\epsilon_{\text{floor}})}, \qquad
  \frac{\partial r}{\partial \phi_x} = -\frac{\partial |\grad
  G|}{\partial\phi_x}.
\end{equation}
Folded into \eqref{eq:gaussnewton} with $\rho = \mathrm{gate}$ (one
instance per incident tet, summed over all $24$; no double-counting
correction is needed since each tet is visited exactly once, from $x$'s own
perspective, unlike the pairwise term in Section~\ref{sec:junction}):
\begin{equation}
  \frac{\partial E_{\text{eik}}}{\partial \phi_x} \mathrel{+}= \mathrm{gate}
  \cdot 2\, r \cdot \frac{\partial r}{\partial\phi_x}, \qquad
  D_{\text{eik}}(x) \mathrel{+}= \mathrm{gate}\cdot 2 \left(\frac{\partial
  r}{\partial\phi_x}\right)^{\!2}.
\end{equation}

\subsection{A known limitation}

$\partial r/\partial\phi_x$ is proportional to $\grad G\cdot w_m$, which
itself shrinks as the confidence field collapses toward the degenerate flat
state ($G_c \to 0$ for all $c$). The restoring force this term applies
therefore weakens as the field approaches exact collapse rather than
staying bounded away from zero: it is a \emph{repeller}, not a hard
barrier. This is an accepted property of the one-sided Gauss--Newton
formulation (a barrier-function alternative was deliberately not used, to
keep the update in \eqref{eq:update} well-defined and finite everywhere,
including exactly at collapse) and is why the Eikonal weight must be kept
meaningfully large relative to the other terms rather than assumed to
"just work" at any positive value.

\subsection{Buffers}

Per tile: the same shared potential/label cache used by every term
(Section~\ref{sec:buffers}) plus one groupshared accumulator pair per
target node for this term's gradient/diagonal contribution. No additional
global buffer is required; the tetrahedron corner tables used to enumerate
the $24$ incident tets are fixed constant tables, not buffers.

\section{Smoothing Term}
\label{sec:smoothing}

\subsection{Purpose}

A discrete-Laplacian-type penalty on local disagreement between a node's
own confidence value and its immediate same-label neighborhood, applied
directly to the confidence field \eqref{eq:confidence} rather than to raw
potentials -- so that disagreement in \emph{label}, not just in potential
magnitude, is treated as a real discontinuity to be smoothed.

\subsection{Stencil}

The stencil combines two distinct neighborhoods, for $26$ neighbors total:
\begin{itemize}
  \item \textbf{face} neighbors ($6$) and \textbf{edge} neighbors ($12$):
  the standard same-sublattice $3\times3\times3$ neighborhood (offset along
  exactly one or exactly two of the three grid axes, respectively) --
  \emph{not} the full $26$-neighbor cube: its $8$ same-sublattice diagonal
  positions (offset along all three axes) are excluded entirely,
  \item \textbf{cross} neighbors ($8$): the node's $8$ nearest neighbors on
  the \emph{opposite} sublattice (the true BCC nearest-neighbor distance) --
  for an A-node these are the $8$ surrounding B cube-centers, for a B-node
  its own $8$ A corners (the same corners \eqref{sec:buffers}'s relabel
  vote already uses),
\end{itemize}
plus the target node itself, for $27$ taps total. Each tap $t$ carries a
fixed scalar weight $k_t$ depending only on its class:
\begin{equation}
  k_{\text{self}} = 192, \qquad k_{\text{face}} = 16 \ (\times 6), \qquad
  k_{\text{edge}} = 8 \ (\times 12), \qquad k_{\text{cross}} = -48\ (\times
  8).
  \label{eq:smoothkernel}
\end{equation}

\subsection{Residual and fold-in}

Because the underlying per-tet-pair confidence-disagreement energy this
kernel reduces from is exactly quadratic in the unknowns, its
Gauss--Newton linearization is exact (not an approximation): define, for
target $x$ with label $\ell$,
\begin{equation}
  \mathrm{Total}(x) \;=\; \sum_{t=0}^{26} k_t\,\sigma(t,\ell)\,\phi_t,
  \label{eq:total}
\end{equation}
the tap potentials $\phi_t$ signed by \eqref{eq:confidence}'s agreement
rule against $x$'s own label $\ell$ (the self tap, $t{=}0$, trivially
agrees with itself and contributes $+192\,\phi_x$). Then directly
\begin{equation}
  \frac{\partial E_{\text{smooth}}}{\partial\phi_x} = \mathrm{Total}(x),
  \qquad D_{\text{smooth}}(x) = k_{\text{self}} = 192,
\end{equation}
with no residual/derivative bookkeeping beyond evaluating
\eqref{eq:total}: the diagonal is the fixed constant $192$ at every node,
independent of the field's current state.

\emph{Consistency check.} For a spatially uniform, single-label field
($\phi_t = \phi$, $\sigma(t,\ell)=+1$ for every tap),
\begin{equation}
  \mathrm{Total} = \phi\left(192 + 6\cdot16 + 12\cdot 8 + 8\cdot(-48)\right)
  = \phi\left(192+96+96-384\right) = 0,
\end{equation}
confirming that a constant same-label field is an exact fixed point of
this term -- it measures only local disagreement from uniformity, as a
proper smoothing/Laplacian-type operator should.

\subsection{Buffers}

Same shared potential/label cache as every other term (Section
\ref{sec:buffers}); one groupshared accumulator per target node holding
$\mathrm{Total}(x)$.

\section{Juncture Straightness Term}
\label{sec:junction}

\subsection{Purpose}

Where three labels meet, the interface should locally look like three flat
sheets meeting along one straight line (the triple junction), not a
curve that kinks from one tetrahedron to the next. This term penalizes
disagreement, between two tetrahedra sharing a genuinely triple-labeled
interface triangle, in their independently-estimated triple-junction
direction.

\subsection{Identifying a junction interface}

Consider two tetrahedra $\tau_A,\tau_B$ sharing a face (three common
corners). Let $l_0 < l_1 < l_2$ be the labels carried by those three shared
corners, canonically sorted ascending -- this instance is only evaluated
if all three labels are pairwise distinct (a genuine triple junction). The
ascending sort is required because the tangent construction below is
sign-sensitive to odd permutations of the label triple, while the
underlying \emph{set} of three labels an interface triangle carries is the
only thing intrinsic to the interface; fixing a canonical order makes the
sign well-defined and consistent regardless of which of the two tets is
treated first.

\subsection{Per-tet triple-junction tangent}

Within one tet $\tau$ (either $\tau_A$ or $\tau_B$), build the three
per-label confidence gradients restricted to $\tau$'s own four corners via
\eqref{eq:shapegrad}:
\begin{equation}
  g_a = \sum_{c\in\tau} G_{l_0}(c)\,w_c(\tau), \quad
  g_b = \sum_{c\in\tau} G_{l_1}(c)\,w_c(\tau), \quad
  g_c = \sum_{c\in\tau} G_{l_2}(c)\,w_c(\tau).
\end{equation}
The two interface planes present in $\tau$ (label $l_0$ vs.\ $l_1$, and
label $l_0$ vs.\ $l_2$) have normals $g_a - g_b$ and $g_a - g_c$; the
triple junction line lies along their common direction,
\begin{equation}
  T_\tau \;=\; (g_a - g_b)\times(g_a - g_c)
  \;=\; g_a\times g_b + g_b\times g_c + g_c\times g_a,
  \label{eq:tangent}
\end{equation}
the second, cyclic-symmetric form following from bilinearity and
antisymmetry of the cross product. $T_\tau$ is invariant under a cyclic
rotation of $(l_0,l_1,l_2)$ and flips sign under a transposition of any
two -- the reason the canonical sort above is required.

\subsection{Residual}

Let $T_A, T_B$ be \eqref{eq:tangent} evaluated in $\tau_A,\tau_B$
respectively (using the same canonical $(l_0,l_1,l_2)$ for both, since they
share the interface). The residual is
\begin{equation}
  C \;=\; T_A \times T_B,
  \label{eq:junctionresidual}
\end{equation}
not $T_A - T_B$: $C$ vanishes exactly when $T_A,T_B$ are parallel
\emph{or} antiparallel (a junction line has no intrinsic "forward"
direction, so antiparallel tangent estimates are equally aligned), and it
does not penalize a pure magnitude difference between the two tets' own
tangent-magnitude estimates, which depends on tet size/shape/potential
scale rather than on whether the junction is actually straight.

\subsection{Derivative and fold-in}

Let $m_A$ (resp.\ $m_B$) be the target's own local corner index within
$\tau_A$ (resp.\ $\tau_B$, only defined if the target is also a corner of
$\tau_B$), and let $s_a,s_b,s_c\in\{+1,-1\}$ be the sign of the target's
own label's agreement with $l_0,l_1,l_2$. Differentiating
\eqref{eq:tangent} with respect to the target's own potential $\phi_x$
(using $\partial g_a/\partial\phi_x = s_a\,w_{m}$, etc., and bilinearity of
the cross product) gives
\begin{equation}
  \frac{\partial T_\tau}{\partial\phi_x} \;=\; w_{m}(\tau) \times
  \Big[(s_c{-}s_b)\,g_a + (s_a{-}s_c)\,g_b + (s_b{-}s_a)\,g_c\Big].
  \label{eq:dtangent}
\end{equation}
Apply \eqref{eq:dtangent} in $\tau_A$ to get $\partial T_A/\partial\phi_x$;
apply it in $\tau_B$ (using $\tau_B$'s own $g_a,g_b,g_c$ and $m_B$) to get
$\partial T_B/\partial\phi_x$ if the target is also a corner of $\tau_B$,
else $\partial T_B/\partial\phi_x = 0$. Then, by the product rule,
\begin{equation}
  \frac{\partial C}{\partial\phi_x} = \frac{\partial T_A}{\partial\phi_x}
  \times T_B \;+\; T_A \times \frac{\partial T_B}{\partial\phi_x}.
\end{equation}
Folded into \eqref{eq:gaussnewton} (with $C\in\R^3$ treated
componentwise, i.e.\ $C\cdot C$ as the scalar residual-squared and
$C\cdot\partial C/\partial\phi_x$ its half-derivative):
\begin{equation}
  \frac{\partial E_{\text{junc}}}{\partial\phi_x} \mathrel{+}= \rho\cdot
  2\left(C\cdot\frac{\partial C}{\partial\phi_x}\right), \qquad
  D_{\text{junc}}(x) \mathrel{+}= \rho\cdot 2\left\lVert
  \frac{\partial C}{\partial\phi_x}\right\rVert^2 .
\end{equation}

\subsection{Stencil and double-counting}

For target $x$, walk its $24$ ring-1 incident tets; for each, test its $4$
face-adjacent-partner relations. A partner tet may itself touch the
target ("ring1--ring1", $36$ distinct such pairs per target) or not
("ring1--ring2", $24$ distinct pairs). A ring1--ring1 pair is discovered
twice in this walk -- once from each member tet's own perspective -- so
its contribution is halved, $\rho = \tfrac12$, to avoid double-counting;
a ring1--ring2 pair is discovered once, $\rho = 1$.

\subsection{Dependency on the Eikonal term}

$C = T_A \times T_B$ vanishes not only when the two tangents are aligned
but also whenever either tangent vanishes outright -- which happens
whenever the underlying per-label gradients $g_a,g_b,g_c$ degenerate (a
collapsing confidence field). This is exactly the degeneracy
Section~\ref{sec:eikonal}'s term exists to resist: keeping $|\grad
G_\ell|$ away from zero at genuine local confidence peaks keeps the $g_a,
g_b,g_c$ vectors that feed this term meaningfully non-zero, so that a
small $C$ reflects real junction alignment rather than a vacuously
satisfied flat field.

\subsection{Buffers}

Reuses the same shared potential/label cache (Section~\ref{sec:buffers});
one groupshared accumulator pair per target node for this term's
gradient/diagonal contribution. No connectivity buffer: tet corners,
incident-tet enumeration, and face-adjacent-partner lookup are all
computed on the fly from fixed constant tables plus the tet/node index.

\section{Connectivity / Volume Term (planned)}
\label{sec:volume}

\subsection{Purpose}

The three terms above constrain local shape (unit gradient, smoothness,
straight junctions) but nothing yet constrains \emph{how much} territory a
label ends up with. In an under-constrained region, all three can be
satisfied while a label's reconstructed volume drifts arbitrarily far from
its intended extent, or a thin connected region pinches off entirely. This
planned fourth term is meant to give each label's region a target
volume (and, eventually, a topological-connectivity constraint) so that
the solve is pulled back toward a sensible territory size rather than
merely a locally-smooth, locally-straight one.

\subsection{Anticipated structure}

Like the terms above, this term is expected to take the same Gauss--Newton
shape: a residual $r_{\text{vol}}$ (or several, per label and per some
local region) built from a per-tet or per-node volume/flux estimate
derived from the same confidence field \eqref{eq:confidence} and the same
per-tet gradient reconstruction \eqref{eq:shapegrad} used above, folded
into the shared accumulators of \eqref{eq:gaussnewton} and combined into
the same update \eqref{eq:update} alongside the three terms already
described. The exact residual definition and stencil are not yet fixed.

\subsection{Dependency on the Eikonal term}

Whatever form the volume/flux estimate takes, it will be computed from
ratios or products of nearby confidence values and/or their gradients. Such
quantities are only numerically meaningful, and only actually reflect real
territory, when the underlying field is non-degenerate -- a volume or flux
estimate built from a field whose gradients have been allowed to collapse
toward zero is ill-conditioned and does not correspond to anything
physical. This term therefore inherits the same dependency on
Section~\ref{sec:eikonal} that Section~\ref{sec:junction} has: the Eikonal
floor must already be holding the field away from collapse before a
connectivity/volume residual can be evaluated meaningfully.

\subsection{Buffers}

Not yet allocated. Expected to need at least one additional groupshared
accumulator pair (gradient/diagonal, mirroring every term above) and,
depending on the final residual's stencil radius, possibly a small
additional per-label scratch quantity computed once per tile and shared
across that tile's targets -- reusing the same tile/halo structure
described in Section~\ref{sec:buffers} rather than introducing a new one.

\section{Shared Buffers and Dispatch Geometry}
\label{sec:buffers}

All four terms operate within the same tile-based dispatch. Each
compute-shader thread group owns a fixed-size \emph{halo}: a $4\times4
\times4$ window of $A$-nodes and a $4\times4\times4$ window of $B$-nodes
(128 nodes total), loaded once per tile from the global node buffers into
groupshared memory before any term runs:

\begin{center}
\begin{tabular}{@{}ll@{}}
\toprule
Groupshared array & Contents \\
\midrule
\texttt{gLabel[128]} & current label of every halo node \\
\texttt{gPot[128]} & current potential of every halo node \\
\bottomrule
\end{tabular}
\end{center}

Each tile is responsible for updating $16$ \emph{target} nodes (the
tile's own $8$ $A$-nodes and $8$ $B$-nodes, strictly interior to the halo
so that every stencil above -- up to the Eikonal/Junction terms' $24$-tet
ring -- stays within the loaded halo without touching global memory
again). Each term accumulates its own gradient/diagonal contribution per
target into its own groupshared pair before the terms are combined per
\eqref{eq:update}:

\begin{center}
\begin{tabular}{@{}ll@{}}
\toprule
Groupshared array pair & Term \\
\midrule
\texttt{gEikonalGrad[16]}, \texttt{gEikonalDiag[16]} & Eikonal floor (Section~\ref{sec:eikonal}) \\
\texttt{gTotal[16]} & Smoothing (Section~\ref{sec:smoothing}; linear, no separate diag array needed) \\
\texttt{gJunctionGrad[16]}, \texttt{gJunctionDiag[16]} & Juncture straightness (Section~\ref{sec:junction}) \\
\emph{(reserved)} & Connectivity/Volume (Section~\ref{sec:volume}) \\
\bottomrule
\end{tabular}
\end{center}

Global buffers (per node, indexed uniformly across both sublattices):
current label, current potential, and their scratch (next-sweep)
counterparts, read once at the start of a tile's work and written once, at
the very end of the sweep, for that tile's $16$ targets.

Scalar parameters, one shared constant buffer read by every term:

\begin{center}
\begin{tabular}{@{}lll@{}}
\toprule
Symbol & Role & Used by \\
\midrule
$w_{\text{smooth}}$ & Smoothing weight & Section~\ref{sec:smoothing} \\
$w_{\text{junc}}$ & Juncture straightness weight & Section~\ref{sec:junction} \\
$w_{\text{eik}}$ & Eikonal floor weight & Section~\ref{sec:eikonal} \\
$w_{\text{vol}}$ \emph{(planned)} & Connectivity/volume weight & Section~\ref{sec:volume} \\
$\epsilon_{\text{floor}}$ & Diagonal floor / gating scale & Sections~\ref{sec:jacobi}, \ref{sec:eikonal} \\
$\Delta_{\max}$ & Per-sweep potential step clamp & Section~\ref{sec:jacobi} \\
\bottomrule
\end{tabular}
\end{center}

\end{document}
